% \usetheme{Frankfurt}
% \usetheme{Madrid}
% \usetheme{Warsaw}
\setbeamercolor{section in head/foot}{fg=black, bg=white}
\definecolor{myblue}{rgb}{0,0,0.75}
\definecolor{light}{rgb}{0.8,0.8,0.8}
\setbeamercolor{frametitle}{fg=myblue, bg=white}
\setbeamertemplate{mini frames}{} % this removes the little circles, since we have too many slides

%the following is to ensure that subsections don't break the navigation bar
\useoutertheme[subsection=false]{miniframes}
\usepackage{etoolbox}
\makeatletter
\patchcmd{\slideentry}{\advance\beamer@xpos by1\relax}{}{}{}
\def\beamer@subsectionentry#1#2#3#4#5{\advance\beamer@xpos by1\relax}%
\makeatother

\usepackage[T1]{fontenc}
\usepackage[latin9]{inputenc}
\usepackage{tikz}  % Used for looping
\usepackage{lipsum} % Used for looping
\usepackage{listings}
\usepackage{booktabs} % nice looking tables
\usepackage{threeparttable} % footnotes for tables
\usepackage{graphicx} % For resizebox
\usepackage{setspace}
\usepackage{epsfig}
\usepackage{comment}
\usepackage{bbold} %for nice indicator function \mathbb{1}
\usepackage{amsmath}
\usepackage{amsfonts}
\usepackage{amssymb}
\usepackage{textcomp}
\usepackage{comment}

% To fix frame count (so the appendix is not counted)
\setbeamercolor{background canvas}{bg=white}
\newcommand{\backupbegin}{
   \newcounter{framenumberappendix}
   \setcounter{framenumberappendix}{\value{framenumber}}
}
\newcommand{\backupend}{
   \addtocounter{framenumberappendix}{-\value{framenumber}}
   \addtocounter{framenumber}{\value{framenumberappendix}} 
}

\setcounter{secnumdepth}{3}
\setcounter{tocdepth}{3}

\makeatletter

\setbeamercovered{transparent}
\setbeamertemplate{footline}[frame number]
\makeatother
\usepackage[english]{babel}
\usepackage{setspace}

\definecolor{light}{rgb}{0.8,0.8,0.8}


%%%%%%%%%%%%%%%%%%%%%%%%%%%%%%%%%%%%%%%%
%%% NEW COMMANDS
%%%%%%%%%%%%%%%%%%%%%%%%%%%%%%%%%%%%%%%%
\newcommand{\bi}{\begin{itemize}}
\newcommand{\ei}{\end{itemize}}
\newcommand{\bnm}{\begin{enumerate}}
\newcommand{\enm}{\end{enumerate}}
\newcommand{\bdm}{\begin{displaymath}}
\newcommand{\edm}{\end{displaymath}}
\newcommand{\bblock}{\begin{block}{}}
\newcommand{\eblock}{\end{block}}
\newcommand{\bcenter}{\begin{center}}
\newcommand{\ecenter}{\end{center}}
\newcommand{\bea}{\begin{eqnarray}}
\newcommand{\eea}{\end{eqnarray}}
\newcommand{\bean}{\begin{eqnarray*}}
\newcommand{\eean}{\end{eqnarray*}}
\newcommand{\redify} {\textcolor[rgb]{0,0,.75}}
\newcommand{\red} {\textcolor[rgb]{1,0,0}}
\newcommand{\blue} {\textcolor[rgb]{0,0,1}}
\newcommand{\green} {\textcolor[rgb]{0,0.55,0}}
\newcommand{\orange} {\textcolor[rgb]{0.8,0.4,0}}
\newcommand{\purple} {\textcolor[rgb]{0.7,0,0.8}}
\newcommand{\gray} {\textcolor{gray}}
\newtheorem*{Nassumption}{Assumption}

\def\argmin{\mathop{\it argmin}}
\def\argmax{\mathop{\it argmax}}

%this allows for the table with columns of fixed width and any alignment
\usepackage{array}
\newcolumntype{L}[1]{>{\raggedright\let\newline\\\arraybackslash\hspace{0pt}}m{#1}}
\newcolumntype{C}[1]{>{\centering\let\newline\\\arraybackslash\hspace{0pt}}m{#1}}
\newcolumntype{R}[1]{>{\raggedleft\let\newline\\\arraybackslash\hspace{0pt}}m{#1}}





